\vfill
\subsection{Vorgehen und Aufbau des Berichts}\label{sec:vorgehen}

Das UI-Redesign der DocSecBox folgte einem nutzerzentrierten Vorgehen in Anlehnung an den Human-Centred-Design-Prozess nach ISO~9241-210 \parencite{iso_9241_210}.

Zu Beginn wurden Projektkontext, Stakeholder-Anforderungen, Zielgruppen, Personas, Wettbewerb und die bestehende Informationsarchitektur analysiert.
Anschließend wurde die bestehende Oberfläche in einem Ist-Usability-Test als Baseline untersucht, um Schwachstellen der Informationsarchitektur, Navigation und Bedienbarkeit empirisch zu erfassen.
Aus diesen Ergebnissen wurden Projekt- und Nutzeranforderungen abgeleitet, in eine überarbeitete Informationsarchitektur sowie ein Variantenkonzept für Desktop und Mobil überführt und über Skizzen, Wireframes und Mockups zu einem lauffähigen High-Fidelity-Prototyp konkretisiert.
Abschließend wurde der Prototyp in einem Soll-Usability-Test mit dem Ist-Zustand verglichen und reflektiert.

Der Bericht folgt dieser Prozesslogik:

\autoref{sec:planung} beschreibt die Planung und Analyse mit Ressourcenplanung, methodischem Vorgehen, Stakeholder-Interview, Zielgruppen- und Persona-Analyse, Wettbewerbsanalyse, Bestandsaufnahme der Informationsarchitektur sowie dem Usability-Test des Ist-Zustands.
\autoref{sec:konzeption} überführt die Analyseergebnisse in das Redesign, einschließlich Soll-Informationsarchitektur, Variantenkonzept, Wireframes, Mockups und visuellem Designsystem.
\autoref{sec:evaluation} bewertet den entwickelten Prototyp anhand eines erneuten Usability-Tests und reflektiert Zielerreichung, Limitationen und nächste Iterationsschritte.
